\documentclass[10pt,a4paper]{amsart}
\usepackage[margin=19mm]{geometry}
\usepackage{amsmath,amssymb,mathtools}
\usepackage[scr=rsfs,cal=euler]{mathalfa}
\usepackage{xfrac}
\usepackage[unicode,hidelinks,bookmarks=false]{hyperref}
\newcommand{\ie}{i.e.\ }
\newcommand{\cf}{cf.\ }
\newcommand{\resp}{respectively\ }
\newcommand{\Subsection}{\subsection}
\providecommand{\nobreakdash}{\nobreak\hbox{-}}
\setlength{\emergencystretch}{2em}
\multlinegap0pt
\usepackage{bm}
\usepackage{bbm}
\usepackage{dsfont}
\usepackage{xspace}
\usepackage{cancel}
\usepackage{mathtools}
\usepackage{amsthm}
\makeatletter
\@ifundefined{thm}{\newtheorem{thm}{Thm}}{}
\@ifundefined{definition}{\newtheorem{definition}[thm]{Definition}}{}
\@ifundefined{lem}{\newtheorem{lem}[thm]{Lemma}}{}
\makeatother
\DeclareMathOperator{\Vol}{\text{Vol}}
\DeclareMathOperator{\defeq}{\overset{\text{\tiny{def}}}{=}}
\DeclareMathOperator{\rad}{\text{rad}}
\title[The inter-universal Teichm\"uller theory and new Diophantine]{The inter-universal Teichm\"uller theory and new Diophantine results over rational numbers. II}
\author{Zhong-Peng Zhou}
\date{Source: arXiv:2510.05448v1 (CC BY 4.0)}
\begin{document}
\maketitle

\noindent\emph{Source and licence.} The inter-universal Teichm\"uller theory and new Diophantine results over rational numbers. II. Version arXiv:2510.05448v1 (CC BY 4.0), distributed under the Creative Commons Attribution 4.0 International licence, as recorded in the arXiv licence metadata for this exact version and shown on the abstract page https://arxiv.org/abs/2510.05448v1. Full rights notice: licenses/arxiv-2510.05448-CC-BY-4.0.txt.\par\smallskip

\noindent\textit{Editorial scope.} Counted unit: the Lem whose source label is \texttt{1-3-lem: local-global inequalities} in the article (source0.tex lines 742--775, source label \texttt{1-3-lem: local-global inequalities}), delivered together with the article's own complete original proof of it (source0.tex lines 776--872, one \texttt{proof} environment of 97 lines). No mathematics is written, repaired or re-derived: statement and proof are the article's own text line for line. Required context retained uncounted: 1-3-def: basic notation for global inequalities (lines 681--738); 1-1-eq1 (lines 695--697). Every definition, display and label of those lines is retained unchanged. Omitted from this package and not redistributed: 1--680, 698--741, 873--2064. Nothing inside a retained excerpt is omitted or altered: every retained body is delivered exactly as the article's lines are written, so no omission receipt of any kind accompanies this package. Citations: the retained excerpts carry no citation command at all, so no bibliography entry of the article is transcribed and no reference is redistributed. Reference adaptation: none was needed and none was made; every reference inside the delivered unit resolves inside it, so no unresolved reference or citation is produced. Printed slips kept literal: no printed word, symbol, hypothesis or number of the counted statement or of its proof was altered. Layout and class: the article's own class is \texttt{article} with packages including subfiles, appendix, footmisc, microtype, tocbibind, stmaryrd, amsmath, verbatim, mathrsfs, mathtools, amsthm, amssymb, faktor, array; the wrapper is the lane's pinned \texttt{amsart} editorial wrapper at 10pt on A4, which restates the theorem-like environments the retained text needs and adds the article's own macro definitions that the retained text needs (\texttt{\textbackslash Vol}, \texttt{\textbackslash defeq}, \texttt{\textbackslash rad}), with extra wrapper packages (bm, bbm, dsfont, xspace, cancel, mathtools). The delivered numbering is the unit's own. Assets: no retained span contains a figure environment, a graphics-inclusion command, a PDF-page inclusion, a table or a TikZ picture; no image file is copied into this package, the package declares no asset dependency and no image file is redistributed with it. Licence: version arXiv:2510.05448v1 of this article is distributed under the Creative Commons Attribution 4.0 International licence, as recorded in the arXiv licence metadata for this exact version and shown on the abstract page https://arxiv.org/abs/2510.05448v1. A full rights notice is delivered at licenses/arxiv-2510.05448-CC-BY-4.0.txt. Characters: every retained excerpt is pure ASCII, so the delivered engine, XeLaTeX with its default Latin Modern text font, reports no missing character.
\begin{definition} \label{1-3-def: basic notation for global inequalities}
Throughout this section, we will use the following notation:

(a) Let $\boldsymbol{N, n_0, u_0, p_N}$ be positive integers, such that $v_p(n_0 \cdot N)\ge u_0$ and $p \ge p_N$ for each $p\mid N$.
For each prime number $l$, write $N_l \defeq \prod_{p:l\mid v_p(N)} p^{v_p(N)}$.

(b) Let $\boldsymbol{S}$ be a finite set of prime numbers $\ge 5$, with cardinality $n \ge 2$; 
let $2 \le \boldsymbol{k} \le n$, $\boldsymbol{n_1(S), n_k(S)}$ be positive integers, $\boldsymbol{S_1}$ be a finite set of prime numbers;
write $p_0$ for the sallest prime number in $S$, and write $k(S)$ for the product of the $k$ smallest prime numbers in $S$. 

Suppose that for each prime $p\in B\setminus S_1$ [for the definition of $B$, cf. (e)], we have $v_p(N) \ge n_1(S)$;
 for any $k$ distinct prime numbers $l_1, ..., l_k \in S$ and each prime $p$, if $l_i \mid v_p(N)$ for $1\le i\le k$, then $v_p(N) \ge n_k(S)$.

(c) Let $\boldsymbol{\lambda} > 0$ be a real number. For each $l\in S$, suppose that we have real numbers $\boldsymbol{\textbf{Vol}(l)} \ge 0$, $\boldsymbol{a_1(l)} \ge \boldsymbol{a_4(l)} > 0$, such that 
\small\begin{align} \label{1-1-eq1}
\frac{1}{\lambda}\log(l^{-v_l(N)} \cdot N/N_l  ) \le  a_1(l) \cdot \log\rad(N) - a_4(l) \cdot \log\rad(N_l) + \Vol(l)
\end{align}\normalsize

\textbf{(Optional)} In the present paper, if not otherwise specified, we shall take $\lambda = 6$; instead of defining $a_1(l)$, $a_4(l)$, we shall fix a positive integer $\boldsymbol{e_0}$, and write 
\begin{gather*}
a_1(l) \defeq \frac{l^2+5l}{l^2+l-12}\cdot  (1-\frac{1}{e_0  l}), \;
a_4(l) \defeq \frac{l^2+5l}{l^2+l-12}\cdot \frac{1}{e_0}(1- \frac{1}{l}) .
\end{gather*}


(d) \textbf{(Optional)} Let $\boldsymbol{u'_0}\ge u_0$ be a positive integer, $\boldsymbol{P}$ be a subset of prime numbers, such that for each $p\in P$, if $p\mid N$, then $v_p(n_0N)\ge u'_0$.
Let $\boldsymbol{n'_1(S)} \ge n_1(S)$ be a positive integer, such that for each $p\in B' \setminus S_1$ [for the definition of $B'$, cf. (e)], we have $v_p(N') = v_p(N) \ge n'_1(S)$.

(e) We shall write
\small\begin{gather*}
A\defeq \{p\in S : p\mid N\}, \;
B\defeq \{p\mid N: \exists l\in S,\, l\mid v_p(N) \}, \;
C\defeq \{p\mid N: p \notin (A\cup B)\}, \\
%
A' \defeq P \cap A, \;
B' \defeq P \cap B, \;
C' \defeq P \cap C, \;
N'\defeq \prod_{p\in P}p^{v_p(N)}, \;
n'_0 \defeq \prod_{p\in P}p^{v_p(n_0)}; \;
N_A \defeq \prod_{p\in A}p^{v_p(N)}, \\
N_B \defeq \prod_{p\in B}p^{v_p(N)},  \;
N_C \defeq \prod_{p\in C}p^{v_p(N)}, \;
N'_A \defeq \prod_{p\in A'}p^{v_p(N')}, \; 
N'_B \defeq \prod_{p\in B'}p^{v_p(N')},  \;
N'_C \defeq \prod_{p\in C'}p^{v_p(N')}; \\
%
a_1 \defeq \frac{\lambda}{n}\sum_{l\in S}a_1(l), \;
a_2 \defeq \frac{\lambda}{n}\sum_{l\in S}\Vol(l), \;
a_3 \defeq \sum_{l\in S}\log(l), \;
a_4 =\lambda\cdot \min_{l\in S}\{ a_4(l) \}, \;
a_5 = \sum_{p\in S_1} \log(p); \\
%
b_1 \defeq \max\{\frac{a_1}{u_0}, \frac{k}{n} + \frac{a_1-a_4}{n_1(S)} \}, \;
b'_1 \defeq \max\{\frac{a_1}{u'_0}, \frac{k}{n} + \frac{a_1-a_4}{n'_1(S)} \} \le b_1, \\
\quad  b_2 \defeq \frac{a_1}{u_0}\cdot \log(n_0) + a_2 + a_1a_3 + (a_1-a_4)a_5 .
\end{gather*}\normalsize

\end{definition}

\begin{lem} \label{1-3-lem: local-global inequalities}
Proceeding with the notation in Definition \ref{1-3-def: basic notation for global inequalities} and suppose that $\log(N) < n_k(S) \cdot \log(p_N)$.

(i) We have
\small\begin{gather*}
\sum_{l\in S}v_l(N)\cdot\log(l) = \log(N_A),
\; \sum_{p\in B} \log(p) \le  \sum_{l\in S} \log\rad(N_l) ,
\;  \sum_{p\in C} \log(p) \le \frac{\log(n_0 N_C)}{u_0}, 
\\ \sum_{l\in S} \log(N_l)\le (k-1)\cdot\log(N_B),
\; \sum_{p\in B} \log(p) \le \frac{1}{n_1(S)} \cdot \log(N_B) + a_5 ,
\\ \sum_{p\in B} \log(p) \le \frac{1}{n_1(S)} \cdot \log(N_B/N'_B) +  \frac{1}{n'_1(S)} \cdot \log(N'_B) + a'_5 .
\end{gather*}\normalsize

(ii) We have
\small\begin{align}  \label{1-1-eq3}
   \log(N) \le a_1\cdot\log\rad(N) - a_4 \sum_{p\in B}\log(p) + a_2 + \frac{k}{n}(\log(N)-\log(N_C)).
\end{align}\normalsize

(iii) We have
\small\begin{align*}
    \log(N) \le a_1\cdot\log\rad(N_C) + (\frac{k}{n}+ \frac{a_1-a_4}{n_1(S)})(\log(N)-\log(N_C))
    + a_2 + a_1a_3 + (a_1-a_4)a_5 .
\end{align*}\normalsize

(iv) We have 
\small\begin{align*}
\log(N) \le &  a_1 \cdot \log\rad(N_C) 
+ (\frac{k}{n} + \frac{a_1-a_4}{n_1(S)})(\log(N/N')-\log(N_C/N'_C))
\\ & + (\frac{k}{n} + \frac{a_1-a_4}{n'_1(S)})(\log(N')-\log(N'_C))
+ a_2 + a_1a_3 + (a_1-a_4)a'_5 .
\end{align*}\normalsize


\end{lem}

\begin{proof}
For assertion (i), by the definition of the set $A$, we have 
\small\begin{align*}
    \sum_{l\in S}v_l(N)\cdot\log(l) = \sum_{l:l\in S,l\mid N}v_l(N)\cdot\log(l) = \sum_{l\in A}v_l(N)\cdot\log(l) = \log(N_A).
\end{align*}\normalsize


By the definition of $N_l$ and $B$, we have 
\small\begin{align*}
\sum_{l\in S} \log\rad(N_l)
= \sum_{l \in S} \sum_{p: p\mid N, l\mid v_p(N)} \log(p)
= \sum_{p\in B}\log(p)\cdot\big(\sum_{l:l \in S,l\mid v_p(N)} 1 \big)
\ge \sum_{p\in B} \log(p) 
\end{align*}\normalsize


By the definition of $N_l$ and $B$, we also have 
\small\begin{align*}
\sum_{l\in S} \log(N_l)
= \sum_{l \in S} \sum_{p:p\mid N,l\mid v_p(N)} v_p(N)\cdot \log(p)
= \sum_{p\in B}v_p(N)\cdot \log(p)\cdot\big(\sum_{l:l \in S,l\mid v_p(N)} 1 \big).
\end{align*}\normalsize


By the definition of $n_0$ and $u_0$, for each $p\in C$, we have $v_p(n_0 \cdot N) \ge u_0$, hence 
\small\begin{align*}
\sum_{p\in C} \log(p) \le&  \frac{1}{u_0} \sum_{p\in C} (v_p(n_0) + v_p(N) ) \cdot \log(p) 
\\ \le&  \frac{1}{u_0} \sum_{p\in C} v_p(N)\cdot \log(p) 
+ \frac{1}{u_0} \sum_{p} v_p(n_0) \cdot \log(p)   
= \frac{\log(n_0 N_C)}{u_0}.
\end{align*}\normalsize


Now for each prime number $p\in B$, we claim that $\sum_{p:p\in S, p\mid v_l(N)} 1 \le k-1$, which can be proved by contradiction as follows.
Assume that $l_1,\dots, l_k\in S$ are $k$ distinct prime numbers, such that $l_i\mid v_l(N)$ for $1\le i\le k$. Then by the definition of $n_k(S)$, we have $v_p(N) \ge n_k(S)$, hence
\begin{align*}
    \log(N) = \log(N) &\ge v_p(N)\cdot \log(p)\ge n_k(S) \cdot \log(p_N)  > \log(N).
\end{align*}
--- a contradiction. Thus the claim is true, hence we have 
\small\begin{align*}
    \sum_{l\in S} \log(N_l)
    &= \sum_{p\in B} v_p(N)\cdot \log(p)\cdot\big(\sum_{l: l \in S, l\mid v_p(N)} 1 \big)
   \\&\le (k-1)\cdot\sum_{p\in B}v_p(N)\cdot \log(p) = (k-1)\cdot\log(N_B).
\end{align*}\normalsize

By the definition of $n_1(S)$ and $n_1(p)$, we have $\log(p) \le \frac{1}{n_1(S)}\cdot v_p(N)\cdot\log(p)$ for each $p\in B\setminus S_1$, and $\frac{n_1(p)}{n_1(S)} \cdot \log(p) \le \frac{1}{n_1(S)}\cdot v_p(N)\cdot\log(p)$ for each $p\in B\cap S_1$. Hence
\small\begin{align*}
\sum_{p\in B} \log(p) \le & \frac{1}{n_1(S)} \cdot \sum_{p\in B} v_p(N)\cdot\log(p) + \sum_{p\in B\cap S_1}  (1-\frac{n_1(p)}{n_1(S)})\cdot\log(p)
 \\ \le & \frac{1}{n_1(S)} \cdot \log(N_B) + a_5 .
\end{align*}\normalsize
Similarly, we have
\small\begin{align*}
\sum_{p\in B} \log(p)   \le & \frac{1}{n_1(S)} \cdot \log(N_B/N'_B) + \frac{1}{n'_1(S)} \cdot \log(N'_B) + a_5 .
\end{align*}\normalsize
This proves assertion (i).

Next, we consider assertion (ii). 
By  (\ref{1-1-eq1}) we have 
\begin{equation} \small
\begin{aligned} \label{1-1-eq2}
\log(N) \le \lambda a_1(l)\cdot \log\rad(N) - \lambda a_4(l) \cdot \log\rad(N_l)  
 + \lambda \Vol(l)  + \log(N_l) + v_l(N) \cdot \log(l) .
\end{aligned}
\end{equation}
Then by taking the average of (\ref{1-1-eq2}) for $l\in S$ and by assertion (i), we have
\begin{equation*} \small
\begin{aligned} 
\log(N) \le &  a_1 \log\rad(N) - a_4  \sum_{l\in S}\log\rad(N_l)  + a_2 
+ \frac{1}{n}\sum_{l\in S}\log(N_l)  + \frac{1}{n}\sum_{l\in S} v_l(N) \cdot \log(l) 
\\ \le &  a_1 \log\rad(N) - a_4  \sum_{p\in B}\log(p)  + a_2 + \frac{1}{n}(\log(N_A)+(k-1)\log(N_B) ).
\end{aligned}
\end{equation*}
Then since $\log(N_A), \log(N_B) \le \log(N) - \log(N_C)$ by the fact $A\cap C = B \cap C =\emptyset$, we have 
\small\begin{align*}
   \log(N) \le a_1 \log\rad(N) - a_4  \sum_{p\in B}\log(p) + a_2 + \frac{k}{n}(\log(N) - \log(N_C)).
\end{align*}\normalsize
This proves assertion (ii).

For assertion (iii), by the definition of $A$, $B$ and $C$, for each $p\mid N$, we have $p\in A\cup B\cup C \subseteq S\cap B \cap C$, hence by (i) we have (note that $a_1\ge a_4$ by definitinon)
\small\begin{align*}
&a_1 \log\rad(N) - a_4  \sum_{p\in B}\log(p) 
\le a_1\sum_{p\in S}\log(p) + (a_1-a_4)\sum_{p\in B}\log(p) + a_1\sum_{p\in C} \log(p)
\\ \le& a_1 a_3 +  (a_1-a_4)(\frac{\log(N_B)}{n_1(S)}  + a_5) + a_1\sum_{p\in C} \log(p).
\end{align*}\normalsize
Then by (\ref{1-1-eq3})  and  $\log(N_B) \le \log(N) - \log(N_C)$   we have  
\begin{equation*} \small
\begin{aligned}
& \log(N) \le   a_1 \log\rad(N) - a_4  \sum_{p\in B}\log(p)  + a_2 + \frac{k}{n}(\log(N) - \log(N_C) )
\\ \le & a_1 a_3 + (a_1-a_4)(\frac{\log(N_B)}{n_1(S)}  + a_5) + a_1\sum_{p\in C} \log(p) + a_2 + \frac{k}{n}(\log(N) - \log(N_C) )
\\ \le& a_1\cdot\log\rad(N_C) + (\frac{k}{n}+ \frac{a_1-a_4}{n_1(S)})(\log(N)-\log(N_C))
+ a_2 + a_1a_3 + (a_1-a_4)a_5 .
\end{aligned}
\end{equation*}
This proves assertion (iii).

The proof of assertion (iv) is similar to the proof of assertion (iii), where we shall make use of $\sum_{p\in B} \log(p) \le \frac{1}{n_1(S)} \cdot \log(N_B/N'_B) +  \frac{1}{n'_1(S)} \cdot \log(N'_B) + a_5$ by (i), and $\log(N_B/N'_B) \le \log(N/N') - \log(N_C/N'_C)$, $\log(N'_B) \le \log(N') - \log(N'_C)$ by the fact $(B\setminus B') \cap (C\setminus C') = B'\cap C' = B \cap C = \emptyset$.
\end{proof}

\end{document}
